% Section du chapitre 4 (ex-annexe C)
% Annexe issue de l'ancien chapitre 3 (restructuration en 3 chapitres)
% ============================================================================
% CHAPITRE 3 — SECTION 3.3 : DÉRIVATION DES VALEURS D'OPPOSABILITÉ
% Fichier : partie3/derivation_opposabilite.tex
% ============================================================================

\section{D\'ERIVATION DES VALEURS D'OPPOSABILIT\'E DES SYST\`EMES DE R\'EF\'ERENCE}
\label{sec:derivation_opposabilite}
% ============================================================================
\begin{sloppypar}

Cette section s'ouvre sur la grandeur dont tout le reste d\'epend. Les
sections~\ref{sec:validation_retrospective} \`a~\ref{sec:reproductibilite}
manipulent trois produits CLO de r\'ef\'erence ($0{,}273$ pour ECDSA,
$0{,}437$ pour Helios, $0{,}572$ pour les zk-SNARK), calcul\'es avec la
fiabilit\'e de v\'erification $\Relver$. Encore
faut-il savoir d'o\`u ils viennent. La pr\'esente section les d\'erive,
syst\`eme par syst\`eme, \`a partir de la d\'efinition m\^eme de l'opposabilit\'e
\'etablie au chapitre~\ref{chap:cadre_formel}.

Helios, les zk-SNARK de type Groth16 et ECDSA, trois syst\`emes cryptographiques
r\'eellement d\'eploy\'es, servent de r\'ef\'erence. Ils ont \'et\'e
retenus parce qu'ils occupent trois r\'egimes distincts du compromis
confidentialit\'e--opposabilit\'e et que leurs sp\'ecifications
publi\'ees autorisent une analyse dimension par dimension. Les valeurs
obtenues sont des estimations \emph{analytiques}, \'etablies \`a partir des
sp\'ecifications de protocole et de contextes institutionnels mod\'elis\'es. Elles ne
reposent sur aucune mesure aupr\`es de juges ou de r\'egulateurs r\'eels~\cite{MINKA2026CLO}
et leur statut est celui \'etabli \`a la
section~\ref{subsubsec:statut_deploiement}.

\begin{boiteArbitre}
L'arbitre remonte \`a la source. Trois chiffres circulent depuis le d\'ebut de
ce chapitre ($0{,}273$, $0{,}437$, $0{,}572$) et il veut savoir d'o\`u ils
sortent. On lui annonce d'embl\'ee qu'il s'agit de d\'erivations, c'est-\`a-dire
d'estimations analytiques, sans relev\'e aupr\`es de juges r\'eels. Il note cette
distinction, qui tiendra pour tout ce qui suit.
\end{boiteArbitre}

% -----------------------------------------------------------------------
\subsection{M\'ethodologie de calcul en trois \'etapes}
\label{sec:methodologie_calcul}
% -----------------------------------------------------------------------

L'opposabilit\'e a \'et\'e d\'efinie \`a la section~\ref{sec:clo} comme la
fraction d'incertitude institutionnelle r\'esolue par la preuve~:
\begin{equation}
\Opp(\mathcal{P}, C) = \frac{I(V;\Sigma|C)}{H(V|C)}
\label{eq:opp_definition_calcul}
\end{equation}
o\`u $V = \mathsf{Interpret}(\Sigma, C)$ d\'esigne le verdict institutionnel,
$\Sigma$ la preuve cryptographique et $C$ le contexte. Le th\'eor\`eme du goulot
d'\'etranglement contextuel, \'etabli au m\^eme endroit, borne le num\'erateur
par $\min\{H(\Sigma|C), H(V|C)\}$ et garantit donc $\Opp \leq 1$, contrainte que chacune des trois d\'erivations ci-dessous doit respecter.
Chaque estimation suit trois \'etapes, appliqu\'ees identiquement aux trois
syst\`emes.

\noindent\textbf{\'Etape 1. Mod\'elisation de l'espace de verdicts.}
\'Enum\'erer les d\'ecisions institutionnelles possibles \`a partir des cadres
l\'egaux et r\'eglementaires pertinents, puis estimer $H(V|C)$ en supposant une
distribution approximativement uniforme sur ces verdicts. Cette hypoth\`ese
d'uniformit\'e est une ligne de base conservatrice en l'absence de donn\'ees sur
la jurisprudence r\'eelle~; ses limites sont discut\'ees \`a la
section~\ref{sec:limites_derivations}.

\noindent\textbf{\'Etape 2. Estimation de l'extraction d'information.}
Analyser ce qu'un interpr\`ete institutionnel, dot\'e d'une connaissance
contextuelle $H(\mathcal{C})$, peut effectivement tirer de la preuve $\Sigma$
compte tenu de son expertise et de ses moyens de v\'erification. L'information
extractible est quantifi\'ee en bits, d\'ecompos\'ee dimension par dimension puis
born\'ee par l'hypoth\`ese d'extraction born\'ee
$I(V;\Sigma|C) \leq I(M;\Sigma|C) + \Delta$ \'etablie au
chapitre~\ref{chap:cadre_formel}.

\noindent\textbf{\'Etape 3. Calcul de l'opposabilit\'e.}
Appliquer la d\'efinition~\eqref{eq:opp_definition_calcul}, puis composer avec
les m\'etriques $\Conf$ et $\Relver$ pour obtenir le produit CRO.

\noindent\textbf{Deux \'echelles d'entropie contextuelle.} Une pr\'ecaution de
notation s'impose avant de calculer. Les valeurs employ\'ees dans les trois
d\'erivations qui suivent ($8$, $11$ et $20$ bits) d\'esignent l'entropie
du \emph{contexte d'esp\`ece}, c'est-\`a-dire la richesse informationnelle effectivement
mobilis\'ee par l'institution consid\'er\'ee face \`a un litige donn\'e. Elles ne
doivent pas \^etre confondues avec les plages du
tableau~\ref{tab:estimations_structurelles_HC}, de $20$ \`a $70$ bits, qui
portent sur l'entropie de la \emph{distribution des contextes} d'un domaine
institutionnel entier et servent au dimensionnement des protocoles de la
section~\ref{sec:protocoles_experimentaux}. Le rapport entre les deux
\'echelles est celui d'un cas particulier \`a la population dont il est tir\'e.
La distinction est sans effet sur les valeurs d'opposabilit\'e calcul\'ees
ci-dessous, o\`u $H(\mathcal{C})$ n'intervient pas dans le quotient, mais elle
conditionne l'interpr\'etation des \'ecarts entre syst\`emes.

% -----------------------------------------------------------------------
\subsection{Helios~: vote \'electronique en contentieux \'electoral}
\label{sec:derivation_helios}
% -----------------------------------------------------------------------

\noindent\textbf{Protocole.} Helios~\cite{Adida08} est un syst\`eme de vote
ouvert reposant sur le chiffrement ElGamal et sur des preuves
z\'ero-connaissance construites sur des groupes de logarithme discret de
$2048$ bits.

\noindent\textbf{Contexte institutionnel mod\'elis\'e.} Tribunal de juridiction
g\'en\'erale statuant sur un litige \'electoral, avec
$H(\mathcal{C}) \approx 8$ bits, soit environ $16$ cat\'egories de preuves
distinguables multipli\'ees par $16$ contextes proc\'eduraux distincts.

\noindent\textbf{Formulation CLO.}
\begin{itemize}
\item Messages~: votes individuels $v \in \{1,\ldots,K\}$
\item Preuves~: $\sigma = (\mathsf{Enc}(v), \pi_{\text{valid}}, \pi_{\text{tally}})$,
      o\`u $\mathsf{Enc}$ d\'esigne un chiffrement ElGamal et $\pi$ des preuves
      z\'ero-connaissance de type Fiat-Shamir
\item Verdicts~: $\mathcal{V} = \{$\'election valide, \'election invalide,
      recomptage ordonn\'e, investigation demand\'ee$\}$
\end{itemize}

\noindent\textbf{\'Etape 1. Estimation de $H(V|C)$.} L'espace de verdicts
comporte quatre d\'ecisions distinctes, d'o\`u
\begin{equation}
H(V|C) = \log_2|\mathcal{V}| = \log_2(4) = 2 \text{ bits}.
\label{eq:hvc_helios}
\end{equation}
Cette t\'etralogie d\'ecisionnelle (accepter les r\'esultats, les annuler,
ordonner un recomptage, exiger une investigation) correspond aux standards
des juridictions \'electorales document\'ees~\cite{MINKA2026CLO}.

\noindent\textbf{\'Etape 2. Estimation de $I(V;\Sigma|C)$.} Un juge de
juridiction g\'en\'erale ne peut v\'erifier ind\'ependamment ni les calculs de
logarithme discret ElGamal sur $2048$ bits, ni les transformations
Fiat-Shamir, ni la correction du d\'ecompte homomorphe. En pratique, la preuve
lui parvient comme un \emph{signal expert binaire} transmis par
t\'emoignage, sous la forme \guillemotleft~l'expert cryptographique confirme la
validit\'e~\guillemotright~contre \guillemotleft~l'expert identifie une faille
cryptographique~\guillemotright. Cette cat\'egorisation binaire transmet
\begin{equation}
I(V;\Sigma|C) \approx 1 \text{ bit}.
\label{eq:ivc_helios}
\end{equation}
Cette pauvret\'e informationnelle d\'ecoule de la conception m\^eme du protocole,
puisque la propri\'et\'e z\'ero-connaissance masque le contenu des votes tout en
attestant la correction du calcul. Le juge ne peut
distinguer ni les nuances proc\'edurales, ni les irr\'egularit\'es partielles, ni
les anomalies de comportement des trustees \`a partir des preuves seules.

\noindent\textbf{\'Etape 3. Calcul de l'opposabilit\'e.}
\begin{equation}
\Opp(\text{Helios},\, \mathcal{C}_{\text{tribunal}})
= \frac{1 \text{ bit}}{2 \text{ bits}} = 0{,}50
\label{eq:opp_helios}
\end{equation}
Les preuves cryptographiques r\'esolvent exactement la moiti\'e de
l'incertitude institutionnelle. Le bit r\'esiduel rel\`eve du raisonnement
juridique sur la conformit\'e proc\'edurale, l'interpr\'etation de l'intention
des \'electeurs et le droit \'electoral de la juridiction. La r\'ealit\'e
institutionnelle de cet \'ecart est document\'ee par les incidents IACR~2025 et
UCL~2009, o\`u des trustees pourtant experts en cryptographie ont rencontr\'e
des difficult\'es op\'erationnelles avec la gestion des cl\'es
distribu\'ees~\cite{MINKA2026CLO}. Le protocole de mesure de la
section~\ref{subsec:protocole_opp} confronterait cette valeur de $0{,}50$ \`a
$\Opp_{\mathrm{emp}}$ sur une \'election r\'eellement d\'eploy\'ee.

\noindent\textbf{Produit CRO.}
\begin{equation}
\Conf \cdot \Relver \cdot \Opp
\approx 0{,}92 \times 0{,}95 \times 0{,}50 \approx 0{,}437
\label{eq:produit_helios}
\end{equation}
La confidentialit\'e $\Conf \approx 0{,}92$ d\'ecoule de la s\'ecurit\'e
s\'emantique d'ElGamal sous l'hypoth\`ese DDH~; la fiabilit\'e
$\Relver \approx 0{,}95$, de la soundness des preuves ZK
($< 2^{-80}$)~\cite{MINKA2026CLO}.

% -----------------------------------------------------------------------
\subsection{zk-SNARK Groth16~: conformit\'e r\'eglementaire}
\label{sec:derivation_zksnark}
% -----------------------------------------------------------------------

\noindent\textbf{Protocole.} Preuves z\'ero-connaissance pour la
v\'erification de conformit\'e pr\'eservant la vie priv\'ee. Une institution
financi\`ere y prouve que l'ensemble de ses transactions satisfait les r\`egles de
lutte contre le blanchiment sans en r\'ev\'eler le d\'etail. Construction
Groth16 sur courbe BN254, preuves de $128$ octets en forme compress\'ee~\cite{Groth16}.

\noindent\textbf{Contexte institutionnel mod\'elis\'e.} Audit r\'eglementaire
par une autorit\'e financi\`ere sp\'ecialis\'ee, avec
$H(\mathcal{C}) \approx 11$ bits, soit environ $64$ cat\'egories de preuves
documentaires multipli\'ees par $32$ cadres proc\'eduraux.

\noindent\textbf{Formulation CLO.}
\begin{itemize}
  \item Messages~: bases de transactions $\{(t_i, a_i, v_i)\}$
        (transactions, comptes, montants)
  \item Preuves~: zk-SNARK $\pi$ (construction Groth16 sur BN254)
  \item Verdicts~:
    \[
      \mathcal{V} = \left\{
      \begin{array}{l}
        \text{Conforme},\ \text{Non-conforme},\
        \text{Investigation requise},\\
        \text{\'Evaluation de p\'enalit\'e},\
        \text{Action corrective},\ \text{Escalade},\\
        \text{Approbation conditionnelle},\
        \text{Audit forensique complet}
      \end{array}
      \right\}
    \]
\end{itemize}

\noindent\textbf{\'Etape 1. Estimation de $H(V|C)$.} Huit d\'ecisions
distinctes, d'o\`u
\begin{equation}
H(V|C) = \log_2(8) = 3 \text{ bits}.
\label{eq:hvc_zksnark}
\end{equation}
Les d\'ecisions r\'eglementaires financi\`eres sont structurellement plus
nuanc\'ees que les verdicts \'electoraux, car un r\'egulateur peut approuver sous
conditions, exiger des mesures correctives sans sanction imm\'ediate ou
escalader vers l'application de la loi. Ces huit cat\'egories correspondent aux
pratiques r\'eglementaires document\'ees.

\noindent\textbf{\'Etape 2. Estimation de $I(V;\Sigma|C)$.} Les op\'erations de
couplage elliptique $e: \mathbb{G}_1 \times \mathbb{G}_2 \to \mathbb{G}_T$, la
validation des param\`etres du trusted setup et la v\'erification des
contraintes R1CS demeurent inaccessibles \`a un r\'egulateur sans formation
sp\'ecialis\'ee. Deux bits restent n\'eanmoins extractibles~:
\begin{enumerate}
\item \textbf{V\'erification computationnelle} ($1$ bit). La preuve passe ou
      \'echoue les contr\^oles de couplage, ce qui fournit une information binaire accessible via
      un v\'erificateur certifi\'e sans comprendre la cryptographie sous-jacente.
\item \textbf{Cat\'egorisation de conformit\'e} ($1$ bit). Il s'agit d'une classification
      grossi\`ere lisible dans la structure de la preuve (\guillemotleft~satisfait toutes les contraintes~\guillemotright,
      \guillemotleft~conformit\'e limite~\guillemotright,
      \guillemotleft~violation claire~\guillemotright), accessible via les
      m\'etadonn\'ees du circuit.
\end{enumerate}
\begin{equation}
I(V;\Sigma|C) \approx 2 \text{ bits}.
\label{eq:ivc_zksnark}
\end{equation}
L'avantage sur Helios ($2$ contre $1$ bit) tient au contexte plus riche
($11$ contre $8$ bits), qui autorise une cat\'egorisation plus fine. L'opacit\'e
cryptographique reste la m\^eme, puisque la garantie z\'ero-connaissance de
Groth16 ne peut \^etre contourn\'ee par l'expertise.

\noindent\textbf{\'Etape 3. Calcul de l'opposabilit\'e.}
\begin{equation}
\Opp(\text{zk-SNARK},\, \mathcal{C}_{\text{r\'egulateur}})
= \frac{2 \text{ bits}}{3 \text{ bits}} \approx 0{,}67
\label{eq:opp_zksnark}
\end{equation}
Le tiers restant exige un jugement r\'eglementaire sur la mat\'erialit\'e, le
risque syst\'emique et la priorit\'e d'application, dimensions inaccessibles
\`a la v\'erification cryptographique.

\noindent\textbf{Produit CRO.}
\begin{equation}
\Conf \cdot \Relver \cdot \Opp
\approx 0{,}89 \times 0{,}96 \times 0{,}67 \approx 0{,}572
\label{eq:produit_zksnark}
\end{equation}
La confidentialit\'e $\Conf \approx 0{,}89$ refl\`ete le z\'ero-connaissance
computationnel sous hypoth\`eses de couplage~; la fiabilit\'e
$\Relver \approx 0{,}96$ est l\'eg\`erement inf\'erieure \`a l'unit\'e en raison des
exigences du trusted setup~\cite{MINKA2026CLO}.

% -----------------------------------------------------------------------
\subsection{ECDSA~: litige contractuel avec infrastructure PKI}
\label{sec:derivation_ecdsa}
% -----------------------------------------------------------------------

\noindent\textbf{Protocole.} Authentification par ECDSA sur courbe NIST
P-256~\cite{Johnson01}.

\noindent\textbf{Contexte institutionnel mod\'elis\'e.} Tribunal civil statuant
sur un litige contractuel, avec infrastructure PKI \'etablie et
$H(\mathcal{C}) \approx 20$ bits, d\'ecompos\'es en compr\'ehension PKI
($\approx 2^{10}$ types de certificats, $10$ bits), pr\'ec\'edents
jurisprudentiels ($\approx 2^{8}$ cat\'egories accumul\'ees en d\'ecennies de
droit des signatures num\'eriques, $8$ bits) et cadres l\'egaux (ESIGN, eIDAS et
r\'eglementations sp\'ecifiques, $2$ bits)~\cite{MINKA2026CLO}.

\noindent\textbf{Formulation CLO.}
\begin{itemize}
\item Messages~: documents contractuels, publics par destination
\item Preuves~: signatures ECDSA $(r, s)$ avec \'equation de v\'erification
      publique $sG = R + eP_{\text{pub}}$ sur NIST P-256
\item Verdicts~: espace d\'ecisionnel multidimensionnel couvrant validit\'e,
      force d'authentification, poids juridique, admissibilit\'e proc\'edurale,
      validit\'e temporelle, confirmation d'autorit\'e et pertinence contextuelle
\end{itemize}

\noindent\textbf{\'Etape 1. Estimation de $H(V|C)$.} L'espace de verdicts est
ici multidimensionnel par nature, le tribunal devant \'evaluer simultan\'ement
plusieurs facettes ind\'ependantes de la preuve num\'erique~:
\begin{equation}
H(V|C) \approx 20 \text{ bits}.
\label{eq:hvc_ecdsa}
\end{equation}
Cette valeur co\"incide avec $H(\mathcal{C}) = 20$ bits, ce qui traduit un
contexte institutionnel calibr\'e pour r\'esoudre exactement le niveau
d'incertitude d\'ecisionnelle qu'il affronte.

\noindent\textbf{\'Etape 2. Estimation de $I(V;\Sigma|C)$.} Contrairement aux
preuves z\'ero-connaissance, ECDSA est \emph{transparent par conception}, car
l'\'equation de v\'erification expose les relations math\'ematiques explicites
entre cl\'e publique, hachage du message et composants de la signature. Le
tribunal en extrait environ $18$ bits r\'epartis sur cinq
dimensions~\cite{MINKA2026CLO}~:
\begin{enumerate}
\item \textbf{V\'erification de propri\'et\'e de la cl\'e publique} ($3$ bits), qui recouvre la
      validation de la cha\^{\i}ne de certificats, \'evaluation de la confiance
      en l'\'emetteur, v\'erification des restrictions d'usage de la cl\'e.
\item \textbf{Liaison message--signature} ($4$ bits), qui recouvre l'int\'egrit\'e du hachage,
      la conformit\'e du format de signature, la validation des param\`etres
      algorithmiques et les propri\'et\'es de non-r\'epudiation.
\item \textbf{Validit\'e temporelle} ($3$ bits), qui recouvre l'horodatage, la p\'eriode de
      validit\'e du certificat et le statut de r\'evocation (CRL/OCSP).
\item \textbf{\'Evaluation de l'autorit\'e l\'egale} ($3$ bits), qui recouvre l'autorisation du
      signataire, les conditions de formation du contrat et la conformit\'e
      juridictionnelle.
\item \textbf{Pertinence contextuelle} ($5$ bits), qui recouvre l'alignement avec le type de
      transaction, les exigences r\'eglementaires, les standards industriels, le profil de
      risque et l'ad\'equation au contexte litigieux.
\end{enumerate}
\begin{equation}
I(V;\Sigma|C) \approx 3 + 4 + 3 + 3 + 5 = 18 \text{ bits}.
\label{eq:ivc_ecdsa}
\end{equation}
ECDSA ne cherche pas \`a masquer l'information sur le signataire et le
message, puisqu'un contrat a pr\'ecis\'ement vocation \`a \^etre interpr\'et\'e par des
tiers. La transparence cryptographique, qui p\'enalise lourdement la
confidentialit\'e, est exactement ce qui rend le syst\`eme institutionnellement
puissant. Il s'agit donc d'un compromis d\'elib\'er\'e de conception.

\noindent\textbf{\'Etape 3. Calcul de l'opposabilit\'e.}
\begin{equation}
\Opp(\text{ECDSA},\, \mathcal{C}_{\text{tribunal+PKI}})
= \frac{18 \text{ bits}}{20 \text{ bits}} = 0{,}90
\label{eq:opp_ecdsa}
\end{equation}
Cette interpr\'etabilit\'e de $90\%$ montre que la s\'ecurit\'e
cryptographique peut coexister avec l'utilit\'e institutionnelle.
Les dix pour cent r\'esiduels refl\`etent la discr\'etion judiciaire
irr\'eductible sur l'interpr\'etation contractuelle, l'\'evaluation de la bonne
foi et les consid\'erations \'equitables, qui transcendent la v\'erification
cryptographique~\cite{MINKA2026CLO}.

\noindent\textbf{Produit CRO.}
\begin{equation}
\Conf \cdot \Relver \cdot \Opp
\approx 0{,}31 \times 0{,}98 \times 0{,}90 \approx 0{,}273
\label{eq:produit_ecdsa}
\end{equation}
La faible confidentialit\'e ($\Conf \approx 0{,}31$) est la cons\'equence
directe de la transparence, puisque les documents contractuels sont publics et que seules
les cl\'es priv\'ees demeurent cach\'ees. La tr\`es haute fiabilit\'e
($\Relver \approx 0{,}98$) tient \`a la v\'erification d\'eterministe, la
contrefa\c{c}on supposant la r\'esolution du probl\`eme ECDLP.

% -----------------------------------------------------------------------
\subsection{\'Etude de cas~: ML-DSA-65 (Dilithium-3) en contexte judiciaire}
\label{sec:etude_cas_dilithium}
% -----------------------------------------------------------------------

\noindent\textbf{Protocole.} Signature post-quantique ML-DSA-65 (FIPS~204,
ex-Dilithium-3, cat\'egorie de s\'ecurit\'e NIST~3) sur r\'eseaux euclidiens module-LWE~\cite{Dilithium2022}.

\noindent\textbf{Contexte institutionnel mod\'elis\'e.} Tribunal civil statuant
sur un litige contractuel avec infrastructure PQC en transition, avec
$H(\mathcal{C}) \approx 10{,}5$ bits \`a l'\'echelle de l'esp\`ece (la table de
r\'ef\'erence retient $20$~bits \`a l'\'echelle du domaine,
tableau~\ref{tab:parametres_reference}).

\noindent\textbf{Formulation CLO.}
\begin{itemize}
\item Messages~: documents contractuels
\item Preuves~: signatures Dilithium-3 avec v\'erification sur r\'eseaux
\item Verdicts~: espace d\'ecisionnel similaire \`a ECDSA, mais avec
      incertitude accrue sur l'interpr\'etation des preuves PQC
\end{itemize}

\noindent\textbf{\'Etape 1. Estimation de $H(V|C)$.} L'espace de verdicts
comporte les m\^emes dimensions qu'ECDSA, mais la granularit\'e est
l\'eg\`erement r\'eduite par l'inadaptation institutionnelle~:
\begin{equation}
H(V|C) \approx 10{,}5 \text{ bits}.
\label{eq:hvc_dilithium}
\end{equation}

\noindent\textbf{\'Etape 2. Estimation de $I(V;\Sigma|C)$.} La v\'erification
Dilithium expose des relations alg\'ebriques sur r\'eseaux, moins famili\`eres
aux acteurs institutionnels que les relations ECDSA. L'extraction est born\'ee
par
\begin{equation}
I(V;\Sigma|C) \approx 8 \text{ bits}.
\label{eq:ivc_dilithium}
\end{equation}

\noindent\textbf{\'Etape 3. Calcul de l'opposabilit\'e.}
\begin{equation}
\Opp^q(\text{Dilithium-3},\, \mathcal{C}_{\text{tribunal PQC}})
= \frac{8 \text{ bits}}{10{,}5 \text{ bits}} \approx 0{,}76
\label{eq:opp_q_dilithium}
\end{equation}

\noindent\textbf{Param\`etres post-quantiques.} Pour ML-DSA-65,
$\lambda = 192$. La signature est une cha\^{\i}ne de bits classique (3\,309 octets),
donc $\eta_q = 0$ (Remarque~\ref{rem:forme_raffinee_pq}). La capacit\'e d'extraction
est fix\'ee, \`a titre illustratif, \`a partir de celle d'ECDSA dans le m\^eme usage
(tableau~\ref{tab:parametres_reference})~:
\begin{align}
\Delta_q &= \Delta + \delta_{\text{enh}} - \delta_{\text{degr}}
\approx 35 - 7 = 28 \text{ bits (hypoth\`ese de travail)}
\label{eq:delta_q_dilithium}\\
\eta_q &= 0
\label{eq:eta_q_dilithium}
\end{align}
La confidentialit\'e et la fiabilit\'e se lisent comme pour ECDSA. Le contrat
sign\'e \'etant public, on a $\Conf^q = 0$ au sens de la
D\'efinition~\ref{def:confidentialite_clo}~; la fiabilit\'e de v\'erification est
quant \`a elle proche de~$1$ (infalsifiabilit\'e sous MLWE/MSIS)~:
\begin{align}
\Conf^q &= 0
\label{eq:conf_q_dilithium}\\
\Relver^q &\approx 0{,}98
\label{eq:rel_q_dilithium}
\end{align}
La formulation initiale de cette \'etude retenait $\Conf^q \approx 0{,}95$, valeur
difficile \`a concilier avec $\Conf \approx 0{,}31$ attribu\'ee \`a ECDSA dans le m\^eme usage.

\noindent\textbf{Lecture.} Avec $\Conf^q = 0$ au sens statistique, le produit
CLO$^q$ est nul et la borne est trivialement respect\'ee. La formulation initiale
calculait $0{,}95 \times 0{,}98 \times 0{,}76 \approx 0{,}71$ et le comparait \`a
$(\Delta_q + \eta_q H(C))/\lambda \approx 0{,}147$~:
\begin{equation}
0{,}71 > 0{,}147.
\label{eq:produit_dilithium}
\end{equation}
Ce d\'epassement s'explique comme pour les syst\`emes classiques. Dans ce sc\'enario,
la signature ML-DSA n'est pas une instance du
Th\'eor\`eme~\ref{thm:trilemme_cro_conditionnel} (message public avec $\Relver$ proche
de~$1$ au lieu de $\Relrec$). La pr\'esente formulation pr\'ecise aussi
l'interpr\'etation initiale, selon laquelle \og la borne ne s'applique qu'\`a la
composante cryptographique\fg{}, puisque le th\'eor\`eme borne le produit entier.
L'\'etude de cas illustre la mani\`ere dont les param\`etres $\Delta_q$ et $\eta_q$
s'estimeraient, sans constituer une v\'erification de la borne.

\begin{boiteArbitre}
L'arbitre note l'\'ecart, puis sa cause. Une signature sur un contrat public
ne cache rien. Sa confidentialit\'e \'etant nulle, la comparer \`a la borne n'a
pas de sens. L'\'etude de cas lui montre comment on estimerait $\Delta_q$,
sans pr\'etendre v\'erifier la borne.
\end{boiteArbitre}

% -----------------------------------------------------------------------
\subsection{Synth\`ese des d\'erivations}
\label{sec:synthese_derivations}
% -----------------------------------------------------------------------

\begin{table}[htbp]
\centering
\caption{Composantes du calcul d'opposabilit\'e pour les trois syst\`emes de
r\'ef\'erence. Toutes les valeurs sont des estimations analytiques \'etablies
\`a partir des sp\'ecifications de protocole et de contextes institutionnels
mod\'elis\'es~\cite{MINKA2026CLO}.}
\label{tab:composantes_opposabilite}
\small
\begin{tabular}{@{}lcccccc@{}}
\toprule
\textbf{Syst\`eme} & $H(\mathcal{C})$ & $H(V|C)$ & $I(V;\Sigma|C)$ & $\Opp$ &
$\Conf$ & \textbf{Produit} \\
\midrule
Helios (tribunal g\'en\'eral) & 8 bits & 2 bits & 1 bit & 0{,}50 & 0{,}92 & 0{,}437 \\
zk-SNARK (r\'egulateur financier) & 11 bits & 3 bits & 2 bits & 0{,}67 & 0{,}89 & 0{,}572 \\
ECDSA (tribunal PKI) & 20 bits & 20 bits & 18 bits & 0{,}90 & 0{,}31 & 0{,}273 \\
\bottomrule
\end{tabular}
\end{table}

\noindent\textbf{Deux \'echelles d'entropie contextuelle.} Les valeurs de
$H(\mathcal{C})$ de ce tableau d\'ecrivent le contexte d'une \emph{d\'ecision
particuli\`ere} (\'echelle de l'esp\`ece)~; celles du chapitre~\ref{chap:cadre_formel}
d\'ecrivent la richesse du contexte institutionnel \emph{g\'en\'eral} (\'echelle du
domaine). Les deux sont r\'eunies dans le tableau~\ref{tab:parametres_reference}.

\noindent\textbf{Mise en regard de deux estimations.} Les produits de la
derni\`ere colonne ont \'et\'e obtenus par la voie \emph{descendante}
(espace de verdicts, information extractible, composition avec $\Conf$ et
$\Relver$). La section~\ref{sec:estimation_params_classiques} propose une
d\'ecomposition additive des capacit\'es $\Delta$ ($35$, $56$ et $73$
bits), dont les ordres de grandeur sont compatibles avec les rapports
$\Delta/\lambda$ correspondants~:
\begin{equation}
\frac{35}{128} \approx 0{,}273, \qquad
\frac{56}{128} \approx 0{,}437, \qquad
\frac{73}{128} \approx 0{,}570.
\label{eq:concordance_voies}
\end{equation}
Cette mise en regard ne vaut pas validation, pour deux raisons. (i)~Les valeurs de $\Delta$ ont
\'et\'e initialement fix\'ees par inversion $\Delta := \lambda\cdot(\Conf\cdot\Relver\cdot\Opp)$
(tableau~\ref{tab:validation_cro_classique}), de sorte que les deux voies ne sont
pas ind\'ependantes. (ii)~Ces syst\`emes ne satisfont pas les hypoth\`eses du
th\'eor\`eme (Remarque~\ref{rem:portee_cas_etude}). En particulier,
$H(V|C) = 2 < \gamma = 8{,}6$ pour Helios et $3 < 15{,}5$ pour les zk-SNARKs, avec
les valeurs de $H(\mathcal{C})$ de cette section. Lus comme des instances, les
produits d\'epasseraient $\Delta/(\gamma\lambda)$ d'un facteur de l'ordre de $8{,}6$,
$15{,}5$ et $19$. La mise en regard montre seulement que les ordres de grandeur
retenus sont mutuellement coh\'erents, ce qui est le r\^ole d'une calibration
illustrative.

% -----------------------------------------------------------------------
\subsection{Limites propres \`a ces d\'erivations}
\label{sec:limites_derivations}
% -----------------------------------------------------------------------

Deux limitations affectent sp\'ecifiquement la m\'ethode de calcul expos\'ee
ci-dessus, en sus de la r\'eserve g\'en\'erale \'etablie \`a la
section~\ref{subsubsec:statut_deploiement}.

\noindent\textbf{Uniformit\'e des verdicts.} L'\'etape~1 suppose une
distribution approximativement uniforme sur l'espace des verdicts, faute de
donn\'ees sur la jurisprudence r\'eelle. La distribution effective des
d\'ecisions judiciaires dans des litiges impliquant des preuves
cryptographiques pourrait \^etre fortement asym\'etrique, comme le sugg\`ere par exemple l'acceptation
quasi syst\'ematique d'une signature ECDSA bien \'etablie. Une
telle asym\'etrie abaisserait $H(V|C)$ et rel\`everait m\'ecaniquement $\Opp$. Les
valeurs report\'ees au tableau~\ref{tab:composantes_opposabilite} en
d\'ependent, dans un sens qui d\'epend aussi de l'effet de l'asym\'etrie sur
$I(V;\Sigma|C)$. Le protocole de la
section~\ref{subsec:protocole_opp} mesure pr\'ecis\'ement la distribution
empirique $H_{\mathrm{emp}}(V|C)$~; lever cette hypoth\`ese fait partie de ses
objectifs.

\noindent\textbf{Caract\`ere statique du mod\`ele.} L'analyse traite
$H(\mathcal{C})$ et $\Delta$ comme des constantes. En r\'ealit\'e, les
institutions apprennent et s'adaptent, car l'accumulation de pr\'ec\'edents
augmente $\Delta$ tandis que le changement technologique modifie $H(\mathcal{C})$. Un
mod\`ele dynamique reste \`a construire~; sa calibration rel\`everait de l'analyse
longitudinale pr\'evue au calendrier (tableau~\ref{tab:calendrier_experimental}),
sur trois \`a cinq ans.

\begin{boiteArbitre}
L'arbitre referme la d\'erivation. Il a vu trois syst\`emes occupant chacun un r\'egime
propre, avec des chiffres d\'efendus dimension par dimension. Il a vu aussi, \`a la fin,
deux limites nomm\'ees sans d\'etour, \`a savoir l'uniformit\'e des verdicts dont
d\'ependent les valeurs ainsi que le caract\`ere statique d'un mod\`ele qui ignore
l'apprentissage des institutions.

Il retient que rien n'est cach\'e, puisque les hypoth\`eses sont pos\'ees, les limites
dites et les protocoles susceptibles de les lever d\'esign\'es. Il y voit une estimation
transparente, dont on conna\^it la port\'ee comme les limites.

Les trois chiffres ont d\'esormais une origine explicite. Ce sont des
calibrations que la suite du chapitre manipule comme telles.
\end{boiteArbitre}
\end{sloppypar}